\section{Finite trials for the pressure}
\label{ctr:section}

The upper branch of the certificate search needs finite descriptions whose
values can be evaluated and which bound the pressure from above.  We define
those descriptions first.  A local three-literal interpolation will prove
their upper bound.  Finite formulas will then supply nearly optimal trials,
although initially with an extra source of randomness shared by the sites.
The final part of this section gives the exact change of law used later to
remove that dependence.

Fix a density $a>0$ and a finite penalty $\beta>0$.  In this section the
formula has $\Poi(an)$ clauses, and the three variable positions of each
clause are independent and uniform in $[n]$, with independent fair signs.
Thus positions may repeat.  If $H_n(\sigma)$ is the number of clauses
violated by $\sigma\in\{0,1\}^n$, put
\begin{equation}
 Z_n=\sum_{\sigma\in\{0,1\}^n}e^{-\beta H_n(\sigma)},
 \qquad L_n=\E\log Z_n,\qquad w=1-e^{-\beta}.
 \label{ctr:pressure-data}
\end{equation}
We will prove that $L_n/n$ has a finite limit $P(a,\beta)$.  The parameter
$w$ lies in $(0,1)$: a violated clause multiplies the weight by
$1-w=e^{-\beta}$.

\subsection{Trial descriptions and their evaluation}

A trial consists of a finite hierarchy of conditional averages and a rule
that produces a pair of numbers in $[0,1]$ at each site.  The site index
distinguishes copies of the sampling generators defined below.  The depth is
an integer $d\ge1$, and its averaging exponents satisfy
\begin{equation}
 0<t_1\le\cdots\le t_d\le1.
 \label{ctr:exponents}
\end{equation}
Thus intermediate trials may have equal exponents or exponent one.  A trial
is called \emph{strict} when $0<t_1<\cdots<t_d<1$.

The sampling data are as follows.  A root variable $z_0$ has an arbitrary
law on a standard Borel space; it stores randomness common to all sites.
At each level $j\in\{1,\ldots,d\}$, a uniform variable $z_j\in[0,1]$ is
also shared by all sites.  A site $v$ has uniform variables
$u_0^v,\ldots,u_d^v\in[0,1]$.  All these uniform generators are mutually
independent and independent of $z_0$.  The same two measurable functions
at every site give the messages
\begin{equation}
 f^\pm(z_0,z_1,\ldots,z_d;u_0^v,u_1^v,\ldots,u_d^v)\in[0,1].
 \label{ctr:messages}
\end{equation}
The two coordinates are indexed by the sign of a literal; they need not
sum to one.  General conditional sampling kernels can be encoded by these
uniform generators, so the resulting messages need not be independent.
Each use of a message pair in a clause will be called a reservoir occurrence.
It receives a fresh site, including when another occurrence is in the same
clause.  Distinct fresh generators may of course produce equal sampled values.

For $t>0$ and a bounded random variable $X$, write
\[
 T_tX=\frac1t\log\E e^{tX}
\]
when the variables integrated by the expectation have been specified.
For a loss depending on finitely many sites, define
\begin{equation}
 \mathcal TX=T_{t_1}\cdots T_{t_d}X.
 \label{ctr:transform}
\end{equation}
The rightmost operation acts first.  At level $j$, it integrates $z_j$ and
all the sites' $u_j^v$, conditional on the preceding levels.  The root
variable, all $u_0^v$, and any further root data are held fixed.  Conditional
expectations make this rule meaningful even when some generators are unused.
For the losses below, the root clause counts give a bound on $X$, so the
same definition applies after conditioning on those counts.

An A-packet has a count $D\sim\Poi(3a)$.  Each of its $D$ clauses contains
one common real Boolean variable $b$, two fresh reservoir occurrences, and
three independent fair signs.  Let $I_h(b)$ indicate that the real literal
of clause $h$ is false at $b$, and let $f_{h1},f_{h2}$ be the messages chosen
by the two reservoir signs.  Its loss is
\begin{equation}
 X_A=\log\sum_{b=0,1}\prod_{h=1}^{D}
       \bigl(1-wI_h(b)f_{h1}f_{h2}\bigr).
 \label{ctr:A}
\end{equation}
A B-packet has a count $D\sim\Poi(2a)$ and three fresh reservoir occurrences
with independent fair signs in every clause.  Its loss is
\begin{equation}
 X_B=\sum_{h=1}^{D}\log\bigl(1-wf_{h1}f_{h2}f_{h3}\bigr).
 \label{ctr:B}
\end{equation}
Counts and signs are root data, independent of the trial generators and
$z_0$; different packets use independent such data.  Each clause factor in
Equations~\eqref{ctr:A}--\eqref{ctr:B} lies in $[e^{-\beta},1]$.  In
particular,
\begin{equation}
 \log2-\beta D\le X_A\le\log2,\qquad
 -\beta D\le X_B\le0,\qquad
 |X_A|,|X_B|\le\log2+\beta D.
 \label{ctr:packet-bound}
\end{equation}
The conventions for $D=0$ give $X_A=\log2$ and $X_B=0$.

For a positive integer $k$, use fresh sites and independent clause data in
each of $k$ packets and define the \emph{total} value
\begin{equation}
 G_k=\E\mathcal T\sum_{i=1}^kX_{A,i}
       -\E\mathcal T\sum_{i=1}^kX_{B,i}.
 \label{ctr:G}
\end{equation}
The outer expectations integrate all root data.  A trial is \emph{site-only}
if its message functions do not use $z_1,\ldots,z_d$.  The shared root
variable $z_0$ remains allowed.  For such a trial, conditional on the root
data, the site variables of distinct packets are independent at each level.
The exponential of a sum therefore factors at every conditional average,
and subsequent root expectation gives
\begin{equation}
 G_k=kG_1\qquad\text{for every site-only trial.}
 \label{ctr:additive}
\end{equation}
For a general trial, $G_k$ remains a total value; the shared nonroot
variables may couple the packets, and we do not use this additivity.

Here is the finite class used by the algorithm.  A \emph{rational trial}
is a strict site-only trial in which $z_0$ is uniform on $[0,1]$, all
exponents are rational, and the message pair is a rational-valued step
function on a finite rational rectangular grid in
\[
 (z_0,u_0,u_1,\ldots,u_d)\in[0,1]^{d+2}.
\]
More explicitly, choose a finite increasing list of rational endpoints
from $0$ to $1$ in each coordinate and give a rational pair in $[0,1]^2$
for every resulting rectangle.  Intervals are half open except at the
right endpoint $1$.  The common coordinate $z_0$ is sampled once and is
shared by every site.  Let $\mathcal Q$ denote this class.  Its finite tables,
grids, and strictly ordered exponent lists have an effective enumeration by
finite strings; their validity is checked by rational comparisons.

\begin{lemma}[Evaluation of a fixed rational trial]
\label{ctr:evaluation}
For rational $a>0$, a positive integer $\beta$, and a finite description
$Q\in\mathcal Q$, the value $G_1(a,\beta;Q)$ is uniformly computable from
these data.  In particular, strict negativity of this value is semidecidable
by certified rational upper bounds.
\end{lemma}

\begin{proof}
Fix the count and signs of one packet.  Every reservoir occurrence uses
its own site coordinates; only the root coordinate is common.  Subdividing
the finitely many grids expresses each conditional integral as a finite sum
with rational weights.  The nested transforms and the remaining root average
are consequently finite expressions using rational arithmetic, exponentials,
and logarithms.  Every exponent is positive, and every clause factor is at
least $e^{-\beta}$.  Thus all logarithms have strictly positive arguments,
and these operations admit uniform rational interval approximations to any
prescribed accuracy.  The finite sign average and the Poisson probabilities
are computable in the same way.

It remains to control the count tail uniformly over the description.  The
conditional transforms preserve the bounds in
Equation~\eqref{ctr:packet-bound}.  For $D\sim\Poi(\lambda)$ and any integer
$M\ge0$,
\begin{align}
 \E\bigl[(\log2+\beta D)\1_{\{D>M\}}\bigr]
 &\le 2^{-M}\E\bigl[(\log2+\beta D)2^D\bigr] \notag\\
 &=2^{-M}e^{\lambda}(\log2+2\beta\lambda),
 \label{ctr:poisson-tail}
\end{align}
where $\lambda=3a$ for the A-packet and $\lambda=2a$ for the B-packet.
Choose $M$ so that the sum of these two computable bounds is smaller than
the desired tail error, and then approximate the finitely many retained terms.
This gives a rational interval containing $G_1$ with arbitrarily small width.
Its upper endpoint eventually becomes negative exactly when $G_1<0$.
\end{proof}

\subsection{A local interpolation bound}

We first establish the pressure limit in the same relaxed clause model.
The proof uses only that a fresh clause chooses its three signed positions
independently.  It also provides the logarithmic expansion used for trial
upper bounds.

\begin{lemma}[Pressure limit]
\label{ctr:pressure}
For every $a>0$ and finite $\beta>0$, the sequence $L_n$ in
Equation~\eqref{ctr:pressure-data} is superadditive.  Hence the finite limit
\begin{equation}
 P(a,\beta)=\lim_{n\to\infty}\frac{L_n}{n}
           =\sup_{n\ge1}\frac{L_n}{n}
 \label{ctr:pressure-limit}
\end{equation}
exists and satisfies $\log2-\beta a\le P(a,\beta)\le\log2$.
\end{lemma}

\begin{proof}
Partition $n=n_1+n_2$ variables into two nonempty blocks and put
$\lambda_i=n_i/n$.  At $s\in[0,1]$, take independent within-block clauses
of means $(1-s)an_i$ and global clauses of mean $san$, and write $Z_s$ for
the resulting partition function.  Brackets denote the
current Gibbs distribution, with independent replicas when several samples
occur.  Adding a clause with violation indicator $I$ changes the log partition
function by
\begin{equation}
 \log(1-w\langle I\rangle)
 =-\sum_{\ell\ge1}\frac{w^\ell}{\ell}\langle I\rangle^\ell.
 \label{ctr:log-series}
\end{equation}
The series is uniformly absolutely summable because $0<w<1$.

For fixed replicas $\sigma^1,\ldots,\sigma^\ell$, let $x_i$ be the fraction
of the $2n_i$ signed literals in block $i$ that are false in every replica.
Then $x_i\in[0,1]$.  A fresh within-block clause has mean replica product
$x_i^3$, whereas a fresh global clause has mean replica product
$(\lambda_1x_1+\lambda_2x_2)^3$.  Poisson differentiation and
Equation~\eqref{ctr:log-series} therefore give
\[
 \frac{\mathrm d}{\mathrm ds}\E\log Z_s
 =an\sum_{\ell\ge1}\frac{w^\ell}{\ell}
   \E\left\langle
     \lambda_1x_1^3+\lambda_2x_2^3
       -(\lambda_1x_1+\lambda_2x_2)^3
   \right\rangle\ge0.
\]
The last inequality is convexity of the cube on $[0,1]$.  Differentiation
is justified because one added clause changes $\log Z_s$ by at most
$\beta$, and the series has the uniform bound just noted.  At $s=0$ the
partition function factors into the two block partition functions; at $s=1$
it has the size-$n$ law.  Thus $L_n\ge L_{n_1}+L_{n_2}$.

If the formula has $M$ clauses, then
$n\log2-\beta M\le\log Z_n\le n\log2$.  Taking expectations gives
$n(\log2-\beta a)\le L_n\le n\log2$.  The superadditive lemma now proves
Equation~\eqref{ctr:pressure-limit}.  The argument specializes sparse
interpolation from \cite{BGT2013}; its nonnegative signed-literal proportions
give the cubic convexity used here.
\end{proof}

\begin{proposition}[Sound trial upper bounds]
\label{ctr:upper}
For every $a>0$ and finite $\beta>0$, each finite-depth strict site-only
trial satisfies
\begin{equation}
 L_n\le nG_1\quad(n\ge1),\qquad P(a,\beta)\le G_1.
 \label{ctr:upper-bound}
\end{equation}
\end{proposition}

We prove this by expressing the nested averages as averages over random
weighted paths.  These weights form a Ruelle probability
cascade (RPC), originating in \cite{Ruelle1987}.  The marked Poisson change of
measure below is the cascade reweighting mechanism of Panchenko and Talagrand
\cite[Lemmas~2.1 and~3.1]{PanchenkoTalagrand2007}.  We give the derivation
because both its conditional marks and its partition prediction rule will
be used again in the structural argument.

\begin{lemma}[Cascade identities]
\label{ctr:rpc}
For $0<t_1<\cdots<t_d<1$, there are random positive weights
$(v_\alpha)_{\alpha\in\N^d}$ with $\sum_\alpha v_\alpha=1$ and the
following properties.

Independent samples from these weights give nested exchangeable partitions
by their common prefixes.  Conditional on the partitions of $n\ge1$ sampled
leaves, the probability that the next sample belongs to a specified already
seen level-$j$ cluster $C$, for $1\le j\le d$, is
\begin{equation}
 \frac{n_C-t_j}{n},
 \label{ctr:prediction}
\end{equation}
where $n_C$ is the number of samples in $C$.  These probabilities, and their
differences between parent and child clusters, determine the partition law.
The ranked cluster masses are the limiting sample frequencies.

Take any root data independent of the cascade weights.  Conditional on
those data, attach independent uniform marks to the nonroot vertices,
independently of the weights; any root marks are included in the root data.
Each mark may be a vector of uniforms.  Let $Y_\alpha$ be the value of the
same bounded measurable function of the path marks and root data.  Then
\begin{equation}
 \E\log\sum_{\alpha\in\N^d}v_\alpha e^{Y_\alpha}
   =\E T_{t_1}\cdots T_{t_d}Y_{\mathrm{path}}.
 \label{ctr:rpc-log}
\end{equation}
This identity holds conditionally on the root data.  It also holds when
$|Y_\alpha|\le B$ for all leaves, where $B$ depends only on the root data,
is finite almost surely, and satisfies $\E B<\infty$.
\end{lemma}

\begin{proof}
At every vertex of depth $j-1$, independently place a Poisson process of
positive child points with intensity $t_j u^{-t_j-1}\,\dd u$.  The
unnormalized weight $W_\alpha$ of a leaf is the product of the points along
its path.  To verify that their total mass is finite, suppose a child subtree
has total mass $S'$ measured from that child.  For $x=t_j$, the mapping
$u\mapsto v=uS'$ turns its child process into one of intensity
\[
 cxv^{-x-1}\,\dd v,\qquad c=\E(S')^x.
\]
The attached subtree law is tilted by $(S')^x/c$ and is independent of
the mapped point.  This is the marked Poisson mapping formula.  At a leaf,
$S'=1$.

If $S$ is the sum of a process with intensity $cxv^{-x-1}\,\dd v$, then
\begin{equation}
 \E e^{-qS}=\exp\{-c\Gamma(1-x)q^x\},\qquad q>0.
 \label{ctr:stable-transform}
\end{equation}
It follows that $0<S<\infty$ almost surely.  For $s\ge1$,
\[
 \Prob(S>s)\le
 \frac{1-\E e^{-S/s}}{1-e^{-1}}\le C_{c,x}s^{-x},
\]
with a finite constant $C_{c,x}$.  For $0<s\le1$,
\[
 \Prob(S\le s)\le e\,\E e^{-S/s}
   =e\exp\{-c\Gamma(1-x)s^{-x}\}.
\]
Consequently $\E S^y<\infty$ for $0\le y<x$, and $\E|\log S|<\infty$.
Since the exponents strictly increase down the tree, induction from the
leaves gives the required moment $\E(S')^{t_j}<\infty$ at every vertex.
The total leaf mass is therefore finite and positive; normalize it to
obtain $v_\alpha=W_\alpha/\sum_\gamma W_\gamma$.

We next derive the prediction rule, including the effect of selection through
a parent.  Consider a node whose child exponent is $x$ and whose parent
exponent is $y<x$, with $y=0$ at the root.  Once this node is selected from
its parent, its point-sum law is tilted by $S^y/\E S^y$.  Its subtree marks
remain independent of the mapped child points, with the individually tilted
laws described above.  Write $(V_i)$ for the mapped points, so that
$S=\sum_iV_i$.  For a specified partition of $n$ samples into $h$ nonempty
children with sizes $n_1,\ldots,n_h$, the probability is
\begin{align}
 &\frac1{\E S^y}
   \E\sum_{i_1,\ldots,i_h\ \mathrm{distinct}}
        S^{y-n}\prod_{j=1}^hV_{i_j}^{n_j} \notag\\
 &\quad=
 \frac{\prod_{j=1}^h\bigl(cx\Gamma(n_j-x)\bigr)}
      {\E S^y\,\Gamma(n-y)}
 \int_0^\infty q^{hx-y-1}
       e^{-c\Gamma(1-x)q^x}\,\dd q.
 \label{ctr:eppf}
\end{align}
Indeed, the Gamma integral represents $S^{-(n-y)}$, and the Poisson
factorial-moment formula evaluates the sum over distinct points.  Increasing
$n_j$ by one changes the displayed expression by the factor
$(n_j-x)/(n-y)$.  The complementary probability of a new child is
$(hx-y)/(n-y)$.  Observations deeper in the tree concern independent subtree
marks and do not change these point predictions.  Along an already seen
chain, multiplication of the within-parent probabilities telescopes to
Equation~\eqref{ctr:prediction}.  In the usual Poisson--Dirichlet notation,
the child partition
seen through its parent has law $\operatorname{PD}(x,-y)$, rather than the
fresh $\operatorname{PD}(x,0)$ law.  The finite prediction probabilities
determine the nested partition laws, and conditional independent sampling
with the strong law recovers their cluster masses as frequencies.

For Equation~\eqref{ctr:rpc-log}, condition on root data and put $y_d=Y$ and
$y_{j-1}=T_{t_j}y_j$ recursively, integrating the level-$j$ marks.  At the
last level, map a Poisson point $u$ to $ue^{y_d}$.  Conditional on the
ancestor marks, the mapped process is the original process scaled by
$e^{y_{d-1}}$, since its intensity is multiplied by
$\E e^{t_dy_d}=e^{t_dy_{d-1}}$.  Discard the integrated marks.  After the
scale is divided out, the remaining descendant total has its unmarked law
and is independent of the retained ancestor marks.  Repeating this mapping
at the preceding levels gives, conditionally on the root data,
\[
 \sum_\alpha W_\alpha e^{Y_\alpha}
 \ \stackrel{\mathrm{law}}{=}\
 e^{T_{t_1}\cdots T_{t_d}Y_{\mathrm{path}}}
       \sum_\alpha W_\alpha.
\]
This is an equality of marginal distributions.  Taking mean logarithms and
subtracting the mean logarithm of the unmarked total proves
Equation~\eqref{ctr:rpc-log}; log integrability was established above.
For a root-conditionally bounded $Y$, apply the identity at each admissible
root value and then integrate its bound.

Children may initially be ordered by raw Poisson points or by total subtree
masses.  These permutations depend only on the unmarked cascade and preserve
ancestry, so independent hierarchical marks have the same conditional law
in either ordering.  A permutation made after multiplication by $e^Y$ need
not preserve the law of retained marks.  The backward calculation discards
exactly the marks it integrates and makes no such independence assertion.
\end{proof}

The comparison uses the diluted interpolation method of Franz and Leone
and of Panchenko and Talagrand~\cite{FranzLeone2003,PanchenkoTalagrand2004},
which extends Guerra's interpolation for Gaussian spin glasses
\cite[arXiv version, Theorems~3 and~5]{Guerra2003}.
Panchenko and Talagrand's general finite-step theorem is stated for even
arity.
For this even-arity class, the companion~\cite[Theorem~2.1 and Corollary~10.2]{OpenAIDilutedPressure2026}
proves the exact finite-hierarchy variational formula under the stated
Panchenko--Talagrand hypotheses, including soft even-$K$ SAT in the Poisson
model with independently sampled variable positions.
Its cavity functional is separate from the three-literal trial functional used here.  For K-SAT, Talagrand observed that the nonnegative literal factors
make the convexity step valid at every clause size
\cite[author manuscript, Section~6.5, Equations~(6.120) and~(6.129)]{Talagrand2011MF1}.
We verify the three-literal calculation in the present trial model,
including the independence required for distinct reservoir occurrences.

\begin{proof}[Proof of Proposition~\ref{ctr:upper}]
Use the cascade of Lemma~\ref{ctr:rpc}.  For each reservoir occurrence,
attach a fresh tree of site marks and a fresh root site uniform.  All these
trees are independent conditional on the root global $z_0$; the site-only
message functions use no nonroot global marks.  At $s\in[0,1]$, take
independent Poisson collections of three kinds of factors:
\[
 \begin{array}{c|c|c}
 \text{kind}&\text{mean count}&\text{factor}\\ \hline
 \text{three real positions}&(1-s)an&1-wI\\
 \text{three reservoir occurrences}&(1-s)2an&1-wf_1f_2f_3\\
 \text{one real position, two reservoir occurrences}&s3an&1-wI_0f_1f_2.
 \end{array}
\]
Real indices are uniform in $[n]$, with independent positions in the first
kind; all signs are independent and fair.  Here $I$ is the violation
indicator of a real clause and $I_0$ indicates that the one real literal
is false.  Let $\mathcal Z_s$ be the sum over real assignments $\sigma$ and
cascade leaves $\alpha$ of $v_\alpha$ times the product of these factors,
using the marks on path $\alpha$ for its messages.

A Gibbs replica for this partition sum consists of a real assignment and
a cascade label.  Conditional on $\ell$ replicas and the existing data,
including $z_0$,
let $x\in[0,1]$ be the fraction of real signed literals false in all their
assignments.  Let $y\in[0,1]$ be the expected product of the $\ell$ signed
messages at one fresh reservoir occurrence evaluated at their labels.
This expectation includes its fresh root uniform, its fresh tree of site
marks, and one fresh fair sign; within the occurrence, the replicas use
the marks on their respective paths.  Distinct occurrences use independent
copies of these data.  Hence the mean replica products for the three rows
of the table are respectively $x^3$, $y^3$, and $xy^2$.

Poisson differentiation and Equation~\eqref{ctr:log-series} now give
\[
 \frac{\mathrm d}{\mathrm ds}\E\log\mathcal Z_s
 =an\sum_{\ell\ge1}\frac{w^\ell}{\ell}
    \E\langle x^3+2y^3-3xy^2\rangle\ge0,
\]
because
\begin{equation}
 x^3+2y^3-3xy^2=(x-y)^2(x+2y)\ge0.
 \label{ctr:cubic}
\end{equation}
Each new factor lies in $[e^{-\beta},1]$, so the same bounded-increment and
absolute-convergence arguments justify the differentiation.  The fresh-site
factorization just used is the one supplied by the site-only class.  We
apply this upper-bound argument only to that class, while later intermediate
trials retain their used nonroot shared variables.

At $s=0$, the real partition sum separates from the cascade sum.  Splitting
the reservoir clause process into $n$ independent B-packets gives
\[
 \E\log\mathcal Z_0
   =L_n+\E\mathcal T\sum_{i=1}^nX_{B,i}
\]
by Equation~\eqref{ctr:rpc-log}.  At $s=1$, splitting by the uniform real
index gives an independent A-packet at each index.  The sum over real bits
then factors before the cascade sum, and the same identity gives
\[
 \E\log\mathcal Z_1=\E\mathcal T\sum_{i=1}^nX_{A,i}.
\]
The nonnegative derivative therefore proves $L_n\le G_n=nG_1$, where the
last equality is Equation~\eqref{ctr:additive}.  Taking the pressure limit
proves the second inequality in Equation~\eqref{ctr:upper-bound}.
\end{proof}

\subsection{Finite reservoirs give nearly optimal values}

We next obtain trials from finite formulas.  The comparison with a finite
reservoir parallels the cavity increment in the variational principle of
Aizenman, Sims, and Starr for a Gaussian class including the
Sherrington--Kirkpatrick model
\cite[arXiv version, Lemma~3]{AizenmanSimsStarr2003}.  We compute the
Poisson clause rates directly.
The construction gives a value per packet at most $P+\epsilon$ for any
chosen $\epsilon>0$, but uses one nonroot variable shared by all sites.
That variable will have to be removed before the trial belongs to
$\mathcal Q$.

\begin{lemma}[Finite-reservoir approximation]
\label{ctr:reservoir}
Fix $a>0$ and finite $\beta>0$.  For each positive integer $k$ and each
$\epsilon>0$, there is a depth-one trial with $t_1=1$ such that
\begin{equation}
 \frac{G_k}{k}\le P(a,\beta)+\epsilon.
 \label{ctr:reservoir-bound}
\end{equation}
Consequently there is a sequence of trials, one for each horizon
$N\to\infty$, with $G_N/N\le P(a,\beta)+o(1)$.  These trials may use a
shared nonroot global variable, and the trial at one horizon need not be
the trial at another.
\end{lemma}

\begin{proof}
Fix positive integers $M,k$ and split a size-$M+k$ formula into $M$ old and
$k$ new variables.  Its all-old clauses have mean
\begin{equation}
 \lambda_{M,k}=a(M+k)\left(\frac{M}{M+k}\right)^3
              =\frac{aM^3}{(M+k)^2}.
 \label{ctr:old-rate}
\end{equation}
Use this old formula as the root disorder $z_0$.  A uniform generator $z_1$
samples an assignment from its Gibbs law and is shared by all sites.  A
reservoir occurrence samples an independent uniform old index using its
root site variable $u_0$, and its two messages are the hard indicators that
the corresponding signed literals are false in that assignment.  Its
level-one site variable may be ignored.  Thus the depth-one transform $T_1$
is exactly a log Gibbs average over the old assignment.  Each occurrence
samples its own old index; accidental repetitions agree with the independent
position law of the formula.

Clauses with exactly one new position have total mean
$3akM^2/(M+k)^2$ and split uniformly among the new indices.  Clauses with
at least two new positions have mean
\[
 \mu_{\ge2}=\frac{ak^2(3M+k)}{(M+k)^2}.
\]
If the latter clauses are omitted, summing over the new bits gives the log
Gibbs average of a product of $k$ one-bit sums.  Replacing the exactly-one-new
mean by $3a$ at each new index produces the $k$ A-packets.  The excess mean
in this replacement is
\[
 \Delta_A=3ak-\frac{3akM^2}{(M+k)^2}
         =\frac{3ak^2(2M+k)}{(M+k)^2}.
\]
On the other hand, a size-$M$ formula is obtained from the same old formula
by adding old clauses of mean $aM-\lambda_{M,k}$.  Replacing this mean by
$2ak$ produces the $k$ B-packets, with excess
\[
 \Delta_B=2ak-(aM-\lambda_{M,k})
         =\frac{ak^2(3M+2k)}{(M+k)^2}.
\]
All these replacements can be coupled by adding independent Poisson
clauses.  An added factor lies in $[e^{-\beta},1]$, so it changes the log of
any of the relevant positive partition sums by at most $\beta$.  At the ideal
rates, the two mean log Gibbs increments from the same old formula are exactly
the A and B terms of $G_k$ for this trial, which depends on $(M,k)$.  Comparing
the original and ideal rates therefore gives
\begin{align}
 \bigl|L_{M+k}-L_M-G_k\bigr|
 &\le\beta(\mu_{\ge2}+\Delta_A+\Delta_B) \notag\\
 &=\frac{a\beta k^2(12M+6k)}{(M+k)^2}
 \le\frac{12a\beta k^2}{M}.
 \label{ctr:reservoir-increment}
\end{align}

For fixed $k$, the pressure limit implies
\[
 \liminf_{M\to\infty}(L_{M+k}-L_M)\le kP(a,\beta).
\]
Indeed, if all sufficiently large increments exceeded $kP$ by one fixed
positive amount, summing along any arithmetic progression of step $k$
would contradict $L_n/n\to P$.  Choose $M$ large along this liminf so that
the increment divided by $k$ is at most $P+\epsilon/2$ and the last term of
Equation~\eqref{ctr:reservoir-increment}, divided by $k$, is at most
$\epsilon/2$.  This proves Equation~\eqref{ctr:reservoir-bound}.  Finally,
apply the result separately at $k=N$ with any error sequence tending to zero.
The reservoir size may depend on $N$ and on that error; no uniform choice is
needed.
\end{proof}

The reservoir trial uses the shared level-one Gibbs assignment and has
$t_1=1$, so it is not yet a rational trial.  The remaining approximation
argument replaces its used nonroot shared variables by independent site
sampling, makes the exponents strict, and then approximates the messages by
finite rational tables.  Theorem~\ref{cc:complete} will establish
\[
 P(a,\beta)=\inf_{Q\in\mathcal Q}G_1(a,\beta;Q).
\]
Together with Proposition~\ref{ctr:upper}, this also gives the same infimum
over all finite-depth strict measurable site-only trials.  The distinction
between soundness of each allowed trial and completeness of the class is
essential to the certificate search.

\subsection{Exact changes of path law}

The next identity compares two packet collections while leaving unused site
generators fresh.  Its recursive densities are the cascade changes of law
studied in \cite[Section~3]{PanchenkoTalagrand2007}.  We prove the identity
directly for all exponents in Equation~\eqref{ctr:exponents}, including
coincident exponents and the value one.

\begin{lemma}[Path posterior]
\label{ctr:posterior}
Let $Y$ use a collection of old sites and let $X$ use fresh sites in one
trial, with both losses bounded conditional on the root data.  Define
\[
 Y_\ell=T_{t_{\ell+1}}\cdots T_{t_d}Y\quad(0\le\ell\le d),
 \qquad Y_d=Y,\quad Y_0=\mathcal TY.
\]
The difference $\mathcal T(Y+X)-\mathcal TY$ is the nested transform of
$X$ when the successive conditional laws at level $\ell$ are changed by
the density
\begin{equation}
 \exp\{t_\ell(Y_\ell-Y_{\ell-1})\}.
 \label{ctr:posterior-density}
\end{equation}
The old site variables may be stored as additional global variables.  In
this representation all fresh site generators retain their independent
uniform laws.  Any property of the original fresh-site kernels and their
conditional uniform site laws that holds for almost every complete global
path survives this change of law.
\end{lemma}

\begin{proof}
Condition on the root data, and let $\mathcal H_\ell$ be the history through
level $\ell$.  By the definition of $Y_{\ell-1}$, the density
in Equation~\eqref{ctr:posterior-density} integrates to one conditional on
the history through level $\ell-1$.  Put
\[
 V_\ell=T_{t_{\ell+1}}\cdots T_{t_d}(Y+X)-Y_\ell,
 \qquad V_d=X.
\]
At each level the identity
\[
 V_{\ell-1}=\frac1{t_\ell}\log\E\left[
   e^{t_\ell(Y_\ell-Y_{\ell-1})}e^{t_\ell V_\ell}
   \,\middle|\,\mathcal H_{\ell-1}\right]
\]
follows by subtracting $Y_{\ell-1}$ from the two logarithms.  Applying it
successively outward proves the asserted nested transform.

Each density uses only globals and old-site variables, because $Y$ does.
It leaves the product uniform law of all fresh site generators unchanged.
The old-site variables can be included in the enlarged global history,
which is a standard Borel random object and can again be sampled by
successive uniform generators.  Its finite-path law is absolutely continuous
with respect to the original global and old-site law.  Hence any original
null set of complete global paths remains null, while the fresh-site kernels
and their conditional uniform laws stay the same.  The argument holds for
almost every root value, as required.
\end{proof}

Write $m=t_1$.  The same operation permits a term $h(z_0,z_1)/m$ at both
endpoints whenever the corresponding exponential normalizer is finite and
positive and the endpoint expressions are integrable.  This term passes
unchanged through every level after the first.  First form the posterior
from $Y$; then tilt its first enlarged global law by $e^h$, normalized
conditional on the root.  Conditional on the preceding history, the later
level laws are unchanged.  Fresh site generators still have their original
laws, because this additional density uses no fresh site variable.

We will use the following consequence to pass from a large horizon to one
fixed packet block.  Partition $N$ packets into complete blocks of $k$ and
a remainder of size less than $k$.  Telescope from the all-B endpoint to
the all-A endpoint, changing one block at a time and including the same
$h/m$ at both endpoints when it is present.  For a complete block, the
remaining packets play the role of $Y$.  Lemma~\ref{ctr:posterior} expresses
the averaged increment as the $G_k$ functional of their posterior trial.
The A and B versions of the fresh block have the same base posterior and
use fresh clause data and fresh sites.  These data are sampled after the
base is fixed; no fresh outcome is part of its conditioning.

For each conditioned base and perturbation, let $\Delta_k^{\mathrm{post}}$
denote this averaged block increment.  The packet bounds give
\begin{equation}
 |\Delta_k^{\mathrm{post}}|
 \le kC_{a,\beta},\qquad C_{a,\beta}=\log2+5a\beta.
 \label{ctr:block-bound}
\end{equation}
Indeed, the transformed A endpoint has absolute value at most
$k\log2+\beta\sum D_{A,i}$ and the transformed B endpoint at most
$\beta\sum D_{B,i}$; the mean counts are
$3ak$ and $2ak$.  The remainder has the same bound with its size, hence at
most $kC_{a,\beta}$.  These statements concern conditional log transforms
and require neither strict exponents nor a cascade representation for the
intermediate trials.

Lemma~\ref{cs:sensitivity} next controls changes of all averaging exponents
without a factor depending on the depth.  Section~\ref{ct:section} then
constructs the pathwise continuation records used by
Lemma~\ref{ct:transition}.  The posterior identity above explains why such
records, once established, remain valid under later changes of the old
packet data.
